% Regression: t2_tcdual — imported current-behavior fixture from the handoff corpus.
% Formula: representation: shifting all colours g_i -> g g_i leaves the physical
% RMP 2011.12127, appendix "The Toric Code and quantum double models",
% eq:app:tcode-rep-dual (ref_ fig_app_tc-dual.pdf).
%
%   Dual quantum-double PEPS: colour variables g in G on the PLAQUETTES;
%   physical spins are differences of adjacent plaquette colours.  The
%   tensor carries virtual colours g_1..g_4 on the four dual sites (drawn
%   as the inner diamond, one squiggle-marked bond per colour) and the four
%   physical legs
%       g_1 g_4^{-1} (west),  g_1 g_2 (depth-up),
%       g_2^{-1} g_3 (east),  g_4 g_3 (depth-down).
%   The representation is G-injective (G-isometric) w.r.t. the regular
%   representation: shifting all colours g_i -> g g_i leaves the physical
%   differences invariant.
%
% Declarative model: a 3x3 oblique sheet keeping only the four mid-side
% sites (the dual plaquette corners); the diamond is the cycle of
% distinguished edges between them; the outer tensor boundary is an
% outline region over the whole block; the four boundary legs are the
% physical colour-difference indices.
\documentclass[border=6pt]{standalone}
\usepackage{tenkz}
\begin{document}
\begin{tenkzlattice}[rows=3, cols=3, frame=oblique, boundary legs]
  % tensor boundary (thin outline in the paper)
  \tnregion[slot=complement, outline]{(1-3,1-3)}
  % keep the four dual sites (1,2),(2,1),(2,3),(3,2); drop the rest
  \tnsite[removed]{(1,1)} \tnsite[removed]{(1,3)}
  \tnsite[removed]{(3,1)} \tnsite[removed]{(3,3)}
  \tnsite[removed]{(2,2)}
  % the dual plaquette: virtual colours g_1 (NW bond), g_2 (NE), g_3 (SE),
  % g_4 (SW) around the diamond
  \tnedge[distinguished]{(1,2)-(2,3)}
  \tnedge[distinguished]{(2,3)-(3,2)}
  \tnedge[distinguished]{(3,2)-(2,1)}
  \tnedge[distinguished]{(2,1)-(1,2)}
\end{tenkzlattice}
\end{document}
